\documentclass[9pt,a4paper,twocolumn,twoside]{rho-class/rho}

% --- 1. Idioma y parches de Babel ---
\usepackage[spanish,es-nodecimaldot,es-noindentfirst]{babel}
\addto\shorthandsspanish{\spanishdeactivate{~<>}}

% --- 2. Parche para evitar el error "Not in outer par mode" ---


% --- 3. Carga de Bibliografía ---
% Asegúrate de que el archivo rho.bib existe en la carpeta
\addbibresource{rho.bib}

% --- 4. Datos del Informe ---

\title{Estudio de la vida media y conversión interna del $Ba^{134} m$}
\author{Edgar Fernández Vigil}
\affil{Física - Facultad de Ciencias - Universidad de Oviedo}
\dates{\today}


\journalname{TEIII - Práctica de Laboratorio}
\journal{Vida media y conversión interna del $Ba^{134} m$ }


% --- 5. Resumen ---
\begin{abstract}
    Resumen del informe
\end{abstract}

\keywords{Física, LaTeX, Biber, Rho-class}

% ----------------------------------------------------------
\begin{document}

    \maketitle
    \thispagestyle{firststyle}

\section{Introducción}
    \rhostart{L}a configuración ha sido exitosa. Ahora puedes citar 
    tus libros usando el comando \cite{PFGPlots} donde 'key' es la 
    etiqueta que tengas en tu archivo \texttt{rho.bib}.

\section{Resultados y discusión}
    Aquí puedes empezar a escribir el contenido de tu práctica.

\section{Conclusiones}
    Discute tus resultados aquí.

% --- 6. Impresión de Bibliografía ---
\newpage

\printbibliography

\end{document}